\begin{textsample}{NEG Dim 5 – global\_north\_2024\_09 – Score 1.00 – global\_north\_2024\_09/114632272.txt}  \label{ex:f5_neg_097}
John Kirby says U.S. working to prevent ' all-out war ' in Middle East amid rising tensions Amid rising tensions in the Middle East, White House National Security Communications Adviser John Kirby said Sunday that the Biden administration is doing"everything we can to try to prevent this from becoming an all-out war there with Hezbollah across that Lebanese border."Israel and Iran-backed Hezbollah traded fire earlier Sunday morning, with an Israeli Defense Forces spokesperson saying that Hezbollah launched 150 rockets toward Israel, reaching deeper into the country than many previous strikes.
In response, the IDF said it was striking"Hezbollah terrorist targets"in Lebanon.
The IDF struck 400 targets on Saturday and said that the attacks will only intensify.
The fresh strikes come as Israeli Prime Minister Benjamin Netanyahu pledges to"take whatever action is necessary to restore security and to bring our people safe back to their homes"near the Lebanese border in the north of the country.
Asked by ABC"This Week"anchor George Stephanopoulos if escalation in the region is inevitable, Kirby @ @ @ @ @ @ @ @ @ @ is still possible.
John Kirby on ABC ' s"This Week."ABC News"We believe that there are better ways to try to get those Israeli citizens back in their homes up in the north, and to keep those that are there, there safely, than a war, than an escalation, then opening up a second front there at that border with Lebanon against Hezbollah,"Kirby said.
But Stephanopoulos pushed back, noting it seems like Netanyahu is not listening to the United States ' s consistent pleas for de-escalation."Look, the prime minister can speak for himself and what -- and what -- what policy he ' s trying to pursue, what operations he ' s trying to conduct.
We ' ll, of course, recognize that the tensions are much higher now than they were even just a few days ago....
But all that does, George, is underscore for us how important it is to try to find a diplomatic solution,"he said. @ @ @ @ @ @ @ @ @ @"to attacks from Israel earlier this week.
In Lebanon and Syria, thousands of people were injured Tuesday by exploding pagers used by Hezbollah members as part of an Israeli operation.
Another round of attacks targeting two-way radios used by the group followed on Wednesday.
The two attacks killed at least 39 people and injured more than 3,000, according to the Lebanese Health Ministry.
Kirby reiterated that the U.S. was"not involved"in these attacks, but refused to say much more than that, saying he would not"get into the details.""I will just say, though, George, that we are watching all of these escalating tensions that have been occurring over the last week or so with great concern, and we want to make sure that we can continue to do everything we can to try to prevent this from becoming an all-out war there with Hezbollah across that Lebanese border,"he said.
A panel of United Nations specialists in international law and human rights has @ @ @ @ @ @ @ @ @ @"booby traps"with the potential of harming civilians.
Israel had a hand in the manufacturing of the devices with this type of"supply chain interdiction"operation having been planned for at least 15 years, a U.S. intelligence source confirmed to ABC News.
In response to a question about the security of U.S. supply chains, Kirby said that President Joe Biden"has made it clear that he wants the American supply chain to be as resilient and as vibrant as possible."The attacks, including Israel ' s Friday strike on a Beirut suburb that took out a top Hezbollah commander, signal a new stage of escalation in the Middle East and raise fears of that they will increase the likelihood of an expanded conflict in the region.
How these recent attacks impact the efforts to achieve a cease-fire between Israeli and terrorist organization Hamas in Gaza remains an open question.
Kirby conceded to Stephanopoulos that,"We are not achieving any progress here in the last week to two weeks, @ @ @ @ @ @ @ @ @ @, doesn't appear to be negotiating in good faith."But it doesn't mean that we ' re not trying,"he added.
Kirby ' s response follows a report from The Wall Street Journal that U.S. officials believe an Israel-Hamas cease-fire deal is unlikely before the end of Biden ' s term.
When asked Friday about the likelihood of a deal, Biden replied,"A lot of things don't look realistic until we get them done."Stephanopoulos also asked Kirby about alleged election meddling efforts by Iran that U.S. security agencies warned about last week.
Kirby said there is"a very robust interagency effort all across the government to deter and to defeat foreign malign actors."The American people ought to know that the federal government is working hand in glove with their local and state officials to ensure the safety and security of their ballots and their election day activities,"Kirby said.

% matched lemmas: 
\end{textsample}
